\section{Využití nástrojů AI}

Vojtěch Matěj:
AI využito: Claude Sonnet 5.5
Využito na:
\begin{itemize}
    \item\textbf{Vytvoření testovacích scénářů:} Vstupních testů pro generování vnitřního kódu.
    \item\textbf{Konzultace nejasností v zadání:} Ověření interpretace pravidel jazyka, gramatických konstrukcí a chování vestavěných funkcí při hraničních vstupech.
\item \textbf{Debugging:} Pomoc při identifikaci příčin paměťových chyb, úniků paměti.
\end{itemize}

\noindent Magdaléna Somorovská: AI využitie: Gemini. Využité na:
\begin{itemize}
    \item \textbf{Konzultácia nejasností v zadaní:} overenie interpretácie pravidiel cieľového jazyka \texttt{IFJcode26}, formátovania konštánt, escape sekvencií a správania jednotlivých inštrukcií.
    \item \textbf{Návrh testovacích scenárov:} konzultácia štruktúry unit testov a overovania výstupu generátora kódu.
    \item \textbf{Debugging:} pomoc pri hľadaní a odstraňovaní chýb v testoch a formátovaní výstupu.
\end{itemize}

\noindent František Fiala:
AI využito: Claude Opus 5.5
Využito ke:
\begin{itemize}
    \item Generování unit testů
    \item Konzultace nejasností v zadání 
    \item Debugging
    \item Pochopení problematiky
\end{itemize}
\noindent Jan Křepinský:
AI model: GPT-6.2.
Použití:
\begin{itemize}
    \item Generování unit a systémových testů
    \item Konzultace nejasností ze zadání
    \item Konzultace pro lepší pochopení problematiky
    \item Debugging
\end{itemize}